\begin{textsample}{NEG Dim 3 – global\_north\_2025\_01 – Score -1.00 – global\_north\_2025\_01/120529979.txt}  \label{ex:f3_neg_211}
Indonesia ' s active role in Gaza : The Jakarta Post As an emerging middle power, Indonesia has the credentials and the track record as a peacemaker.
Now under President Prabowo, we have a leader who will not shy away from taking on greater global responsibility.
We should put our"independent and active"foreign policy to the test.
The Jakarta Post A woman sits by rubble as she waits along Gaza ' s coastal al-Rashid Street for people to cross from the Israeli-blocked Netzarim corridor from the southern Gaza Strip into Gaza City on January 26, 2025.
PHOTO : AFP January 27, 2025 JAKARTA -- The ceasefire in Gaza, as precarious as it may seem, presents a rare opportunity for Indonesia to play a bigger role in the search for a lasting peaceful solution to the Israeli-Palestinian conflict, more than what we have achieved so far, which has largely been confined to calling for and endorsing resolutions in the United Nations.
We are talking not only about sending relief supplies @ @ @ @ @ @ @ @ @ @ millions of Palestinians displaced by the Gaza war, and we are talking not only about offering to send \textbf{peacekeepers} under a UN mandate.
We should already be doing all of that, and we are, to some extent.
We are talking about Indonesia beginning to carve out a role in the peace-making process itself.
As an emerging middle power, Indonesia has the credentials and the track record as a peacemaker.
Now under President Prabowo Subianto, we have a leader who will not shy away from taking on greater global responsibility.
We should put our"independent and active"foreign policy to the test.
Obviously, the priorities for the coming weeks are for both sides of the conflict, meaning the Israeli government and Hamas, to implement the terms of the truce agreement, including the return of hostages and prisoners, letting food supplies and other aid into Gaza, and allowing displaced Palestinians to return to their homes.
Now that there are strong pressures for both sides to abide by the terms @ @ @ @ @ @ @ @ @ @ what should follow beyond the truce.
There is the massive reconstruction of Gaza and the healing process on both sides of the war, but beyond this, we should be looking at the efforts at resolving the conflict once and for all.
For now, the two-state proposal remains the only viable solution even though many Israelis and Palestinians are losing faith in the idea.
Now is a good time for Indonesia to assert itself in the peace process.
The idea of Indonesia temporarily accepting some of the displaced Palestinians in Gaza while reconstruction takes place is not as preposterous as some politicians make it out to be.
Indonesia has been mentioned as one of the countries that are being considered for the temporary relocation by Steve Witkoff, the United States special envoy on the Middle East.
This prompted strong reactions in Indonesia and the Arab world, where it is seen as a cynical attempt to force Palestinians off their land.
Putting aside politics, there is a compelling humanitarian reason to relocate some, though @ @ @ @ @ @ @ @ @ @ Prabowo, as president-elect in June already suggested that Indonesia was prepared to take on 1,000 Palestinians who needed medical treatment, and another 1,000 traumatized Palestinian children to be treated and given an education in Indonesia.
Prabowo had conveyed the idea to the Palestinian leaders he met at an international conference on Gaza in Amman.
If that offer still stands, Indonesia needs to reiterate it.
This will give an entry point for Indonesia to get more involved in the peacemaking efforts.
Other big players have tried and failed for the last nearly 80 years.
The peace process in the Middle East needs fresh ideas and initiatives and there may be new players too.
Indonesia is now well placed to take up that role.
Relocating some Palestinians temporarily would be a gesture to show that Indonesia is willing to put its money where its mouth is, as it seeks a more active role in the peace process.
Indonesia can learn from Qatar, which negotiated the truce agreement, that it can be a mediator even in @ @ @ @ @ @ @ @ @ @ there are already informal contacts between Indonesia and Israel and Hamas, so it would not be that difficult for Indonesia to start communication with both sides.

% matched lemmas: peacekeeper
\end{textsample}
